% =====================================================================
%  MSCA PAPERS MAP -- which CERTIFLIGHT concept becomes which paper.
%
%  Standalone. Set as Overleaf's main document; pdfLaTeX + BibTeX.
% =====================================================================
\documentclass[11pt,a4paper]{article}
% =====================================================================
%  Shared preamble for the three research proposals.
%
%  Each proposal \input{}s this file and is otherwise standalone, so any
%  one of them can be set as Overleaf's main document and compiled on its
%  own. Nothing here is specific to a single proposal.
% =====================================================================
\usepackage[T1]{fontenc}
\usepackage[utf8]{inputenc}
% Times, matching the parent manuscript. mathptmx is in every TeX Live and on
% Overleaf; lmodern is deliberately not used, since it is absent from minimal
% distributions and this preamble should compile anywhere.
\usepackage{mathptmx}
\usepackage[a4paper,margin=25mm]{geometry}
\usepackage{amsmath,amssymb,amsthm}
\usepackage{booktabs}
\usepackage{tabularx}
\usepackage{xcolor}
\usepackage{titlesec}
\usepackage{enumitem}
\usepackage{microtype}
\usepackage[hidelinks,breaklinks]{hyperref}

% ---- shared style tokens, matching the parent manuscript ------------
\definecolor{netcol}{RGB}{31,78,121}
\definecolor{autocol}{RGB}{18,120,90}
\definecolor{bridgecol}{RGB}{176,96,0}
\definecolor{accent}{RGB}{140,20,20}

\newcommand{\layernet}{\textcolor{netcol}{\textbf{network}}}
\newcommand{\layerauto}{\textcolor{autocol}{\textbf{autonomy}}}
\newcommand{\dataset}[1]{\textsc{#1}}
\newcommand{\code}[1]{\texttt{#1}}
\newcommand{\Real}{\mathbb{R}}
\newcommand{\Norm}[1]{\left\lVert #1\right\rVert}
\newcommand{\MCR}{\mathrm{MCR}}

\newtheorem{claim}{Claim}
\newtheorem{proposition}{Proposition}
\newtheorem{definition}{Definition}
\theoremstyle{remark}
\newtheorem{remark}{Remark}

% ---- section styling -------------------------------------------------
\titleformat{\section}{\large\bfseries\color{netcol}}{\thesection}{0.7em}{}
\titleformat{\subsection}{\normalsize\bfseries}{\thesubsection}{0.6em}{}
\titlespacing*{\section}{0pt}{1.4em}{0.6em}

\setlist[itemize]{leftmargin=1.35em,topsep=0.3em,itemsep=0.25em}
\setlist[enumerate]{leftmargin=1.6em,topsep=0.3em,itemsep=0.25em}

% ---- a boxed "what this rests on" callout ---------------------------
\newenvironment{honestly}
  {\par\medskip\noindent\begin{tabular}{@{}p{2.5mm}@{\hspace{3mm}}p{0.9\linewidth}@{}}
   \textcolor{accent}{\rule{2.5mm}{\baselineskip}} &
   \small\itshape}
  {\end{tabular}\par\medskip}

% ---- the proposal header --------------------------------------------
\newcommand{\proposalheader}[3]{%
  \begin{center}
  {\Large\bfseries #1\par}\medskip
  {\large\itshape #2\par}\medskip
  {\small #3\par}
  \end{center}
  \medskip\hrule\medskip}

\usepackage{array}
\newcolumntype{Y}{>{\raggedright\arraybackslash}X}
\newcolumntype{P}[1]{>{\raggedright\arraybackslash}p{#1}}

\begin{document}

\proposalheader
  {MSCA $\to$ Papers}
  {Two Conference Papers from the CERTIFLIGHT Proposal}
  {RobustUAVs.ai research programme $\cdot$ Roger Nick Anaedevha $\cdot$
   2026-10-01\\
   Companion to the integration plan for Papers~I and II}

% =====================================================================
\section{The decision}
% =====================================================================
The CERTIFLIGHT MSCA proposal has four objectives (O1--O4) resting on four
theorems (T1--T4) across five work packages. Two of them carry a self-contained
technical result strong enough for a top venue \emph{now}, on the preliminary
work already in hand. Each becomes one paper.

\begin{itemize}
\item \textbf{Paper~III} (NeurIPS) --- from O3/T3/T4, WP4.
  \emph{Certificates as Constraints.} The certified-constrained multi-agent
  controller: a formal robustness certificate used as a time-varying, pointwise
  CMDP constraint, with feasibility preserved during learning.
\item \textbf{Paper~IV} (USENIX Security) --- from O1/O2, T1/T2, WP2/WP3.
  \emph{When the Attacker Is Indistinguishable from the Weather.} The unified
  perturbation budget and cross-layer benchmark, reframed around the security
  observation that an attacker can disguise interference as a fault or a gust.
\end{itemize}

\noindent
O4/WP5 (the assurance mapping) is \emph{not} a third paper here: it is
Paper~II of the earlier batch (the SoK), which already absorbed it. T1's
tightness-gap result lives in Paper~IV; T2's interface-rate result is already in
the parent manuscript and in Paper~I.

% =====================================================================
\section{Why these two venues}
% =====================================================================
\begin{center}\small
\begin{tabularx}{\linewidth}{@{}P{0.12\linewidth}P{0.16\linewidth}Y@{}}
\toprule
\textbf{Paper} & \textbf{Venue} & \textbf{Why it fits} \\
\midrule
III & NeurIPS & The result is about the \emph{quantifier} on a learned
constraint --- pointwise vs.\ in-expectation, under a constraint set that is
itself a certified, adversarially moving object. A safe-RL / learning-theory
contribution. The swarm is the setting, not the claim. \\
IV & USENIX Security & A threat-model correction (the attacker who dresses as the
weather), an evaluation methodology with enforced provenance, and a released
benchmark artifact. The shape USENIX rewards, and a venue where an evaluation
substrate can become a community default. \\
\bottomrule
\end{tabularx}
\end{center}

% =====================================================================
\section{Which MSCA theorem each paper carries}
% =====================================================================
\begin{center}\small
\begin{tabularx}{\linewidth}{@{}P{0.10\linewidth}P{0.40\linewidth}Y@{}}
\toprule
\textbf{MSCA} & \textbf{What it is} & \textbf{Where it lands} \\
\midrule
T1 & Heterogeneous-source composition with a quantified tightness gap &
\textbf{Paper~IV}, as the Makarov / VaR-aggregation characterisation of the
union-bound slack. \\
T2 & Interface theorem, discrete detection $\to$ continuous budget;
$\gamma$ as a supremum over an ODD & Parent manuscript (proved) and
\textbf{Paper~I} (the $\gamma$ side); the budget side is \textbf{Paper~IV}. \\
T3 & Feasibility preservation for certified-constrained policies &
\textbf{Paper~III}, Propositions~1--2. \\
T4 & Coverage-degradation bound for fleets with degraded agents &
\textbf{Paper~III}, Propositions~3--4. \\
\bottomrule
\end{tabularx}
\end{center}

% =====================================================================
\section{How the four papers relate}
% =====================================================================
\begin{center}\small
\begin{tabularx}{\linewidth}{@{}P{0.12\linewidth}Y@{}}
\toprule
\textbf{Paper} & \textbf{Object, and its place in the chain
$\theta\to\Delta\to\delta\to\mathrm{MCR}$} \\
\midrule
Parent & Proves the chain for the time-delay class. \\
I (S\&P) & Treats $\Delta$ and $\gamma$ as \emph{design variables}
(authentication as credit and debit). \\
II (SoK) & Systematises the whole field by what its guarantees
\emph{quantify over}. \\
IV (USENIX) & Broadens the \emph{source} of $\delta$ (fault, gust, disguised
attack) and builds the benchmark that measures it. \\
III (NeurIPS) & Consumes the per-agent radius the chain produces, as a
\emph{fleet controller} constraint. \\
\bottomrule
\end{tabularx}
\end{center}

\noindent
The four are non-overlapping in object: a certificate with fixed inputs
(parent), a certificate with tunable inputs (I), the field's certificates in
general (II), a certificate against disguised sources (IV), and a controller
that consumes certificates (III). Each cites the others; none repeats another's
result.

% =====================================================================
\section{Readiness, honestly}
% =====================================================================
\begin{center}\small
\begin{tabularx}{\linewidth}{@{}P{0.10\linewidth}Y Y@{}}
\toprule
\textbf{Paper} & \textbf{Done} & \textbf{To do before submission} \\
\midrule
III & Propositions 1--4 proved and brute-force-checked (3000/3000);
simulator and certificate exist &
the CMDP wrapper, projection layer, degradation schedule, and the RL study \\
IV & Six-corpus ingest, certificate engine, residual-budget lemma, provenance
flags, $\gamma$ campaign &
the fault/environment budget terms, the Makarov reduction (or its empirical
fallback), the evasion construction, benchmark packaging and release \\
\bottomrule
\end{tabularx}
\end{center}

\begin{honestly}
Two claims are not yet established and are flagged as such in the papers. In
Paper~IV, that the three-source composition reduces cleanly to a
fixed-marginals-unknown-copula problem (Makarov) is a reduction to prove, with
the sound union bound as the stated fallback. In Paper~III, the empirical RL
results do not yet exist; the three propositions are the contribution and already
hold, so the paper degrades to a theory-plus-feasibility-layer result if MARL
training is unstable at the target fleet size. Neither paper's headline depends
on an unproven claim.
\end{honestly}

\end{document}
